\documentclass[11pt,a4paper]{article}

\usepackage[utf8]{inputenc}
\usepackage[T1]{fontenc}
\usepackage[margin=2.5cm]{geometry}
\usepackage{hyperref}
\usepackage{enumitem}
\usepackage{xcolor}
\usepackage{titlesec}

\hypersetup{
  colorlinks=true,
  linkcolor=blue!50!black,
  urlcolor=blue!50!black,
  pdftitle={Guia de Uso - YT Music Sync},
  pdfauthor={YT Music Sync}
}

\titleformat{\section}{\Large\bfseries}{\thesection}{1em}{}
\titleformat{\subsection}{\large\bfseries}{\thesubsection}{1em}{}

\newcommand{\dica}[1]{%
  \par\medskip\noindent
  \fcolorbox{blue!40!black}{blue!5}{%
    \begin{minipage}{\dimexpr\linewidth-2\fboxsep-2\fboxrule}
      \textbf{Dica:} #1
    \end{minipage}%
  }\par\medskip
}

\newcommand{\atencao}[1]{%
  \par\medskip\noindent
  \fcolorbox{orange!60!black}{orange!8}{%
    \begin{minipage}{\dimexpr\linewidth-2\fboxsep-2\fboxrule}
      \textbf{Atenção:} #1
    \end{minipage}%
  }\par\medskip
}

\title{\textbf{Guia de Uso --- YT Music Sync}\\[0.3em]
\large Como duas pessoas sincronizam o YouTube Music em tempo real}
\author{}
\date{\today}

\begin{document}

\maketitle

\section*{Visão geral}

O YT Music Sync é uma extensão de navegador que sincroniza a reprodução do
YouTube Music (play, pause e seek) entre duas pessoas em tempo real. Quando
uma pessoa dá play, pause ou pula pra outro trecho da música, o mesmo
acontece automaticamente no player da outra pessoa.

O sistema tem duas partes:

\begin{itemize}[nosep]
  \item \textbf{A extensão}, instalada no navegador de cada uma das duas
        pessoas (Chrome ou Firefox).
  \item \textbf{O backend} (um Cloudflare Worker), que existe uma única vez
        e serve de ponte entre as duas extensões.
\end{itemize}

\dica{Só é preciso fazer o deploy do backend \textbf{uma vez}. Depois disso,
as duas pessoas só precisam instalar e configurar a extensão.}

\section{Pré-requisitos}

\begin{itemize}
  \item Uma conta na \href{https://dash.cloudflare.com/sign-up}{Cloudflare}
        (gratuita, sem cartão de crédito).
  \item Node.js instalado --- só na máquina de quem for fazer o deploy do
        backend (Passo 1).
  \item Google Chrome e/ou Mozilla Firefox instalado nas duas máquinas.
  \item O código-fonte do projeto (repositório
        \texttt{sync-extension-ytbmusic}).
\end{itemize}

\section{Passo 1 --- Publicar o backend (feito uma única vez)}

Uma das duas pessoas (só uma, não precisa repetir) executa, dentro da pasta
\texttt{worker} do projeto:

\begin{verbatim}
cd worker
npm install
npx wrangler login
npx wrangler deploy
\end{verbatim}

O comando \texttt{wrangler login} abre o navegador pedindo autorização na
conta Cloudflare. O \texttt{wrangler deploy} publica o backend e imprime uma
URL parecida com:

\begin{verbatim}
https://ytmusicsync-worker.SEU-SUBDOMINIO.workers.dev
\end{verbatim}

\textbf{Anote essa URL trocando \texttt{https://} por \texttt{wss://}} --- é
o endereço que as duas pessoas vão usar na extensão:

\begin{verbatim}
wss://ytmusicsync-worker.SEU-SUBDOMINIO.workers.dev
\end{verbatim}

\section{Passo 2 --- Combinar um código de sala}

As duas pessoas precisam concordar em um \emph{código de sala}: qualquer
palavra ou frase, igual para os dois (por exemplo \texttt{sala-do-joao-e-maria}).
Esse código funciona como uma senha simples --- só quem souber o código e a
URL do backend consegue entrar na mesma sala de sincronização.

\atencao{Não existe login nem autenticação de usuário no sistema (por
decisão de projeto). O código de sala é o único controle de acesso, então
escolham algo que não seja óbvio de adivinhar.}

\section{Passo 3 --- Instalar a extensão}

Este passo é repetido pelas \textbf{duas pessoas}, cada uma no seu próprio
navegador.

\subsection{No Google Chrome}

\begin{enumerate}
  \item Abrir \texttt{chrome://extensions}.
  \item Ativar o interruptor \textbf{``Modo do desenvolvedor''} no canto
        superior direito.
  \item Clicar em \textbf{``Carregar sem compactação''}.
  \item Selecionar a pasta \texttt{extension} do projeto.
\end{enumerate}

\subsection{No Mozilla Firefox}

\begin{enumerate}
  \item Abrir \texttt{about:debugging\#/runtime/this-firefox}.
  \item Clicar em \textbf{``Carregar extensão temporária''}.
  \item Selecionar o arquivo \texttt{manifest.json} dentro da pasta
        \texttt{extension}.
\end{enumerate}

\atencao{No Firefox, uma extensão carregada como ``temporária'' é removida
sempre que o navegador é fechado. É preciso repetir esses 3 passos toda vez
que o Firefox reiniciar, a não ser que a extensão seja assinada e instalada
de forma permanente.}

\section{Passo 4 --- Configurar a extensão}

Ainda nas \textbf{duas} instalações, cada pessoa deve:

\begin{enumerate}
  \item Clicar no ícone da extensão \textbf{YT Music Sync} na barra do
        navegador.
  \item Preencher o campo \textbf{Worker URL} com a URL \texttt{wss://...}
        anotada no Passo 1.
  \item Preencher o campo \textbf{Código da sala} com o código combinado no
        Passo 2.
  \item Clicar em \textbf{Salvar}.
\end{enumerate}

\dica{Os valores de \textbf{Worker URL} e \textbf{Código da sala} precisam
ser \emph{idênticos} nas duas instalações. Um espaço a mais ou uma letra
diferente já impede a sincronização.}

\section{Passo 5 --- Usar}

Com a extensão instalada e configurada nos dois lados:

\begin{enumerate}
  \item As duas pessoas abrem \url{https://music.youtube.com}.
  \item Uma delas toca uma música, pausa, ou pula pra outro trecho.
  \item O mesmo comando aparece automaticamente no player da outra pessoa.
\end{enumerate}

Não importa quem inicia a reprodução --- a sincronização funciona nos dois
sentidos, a qualquer momento.

\section{Limitações importantes}

\begin{itemize}
  \item Funciona apenas com \textbf{duas pessoas} por sala.
  \item Não guarda histórico de sessões --- o estado da sala existe apenas
        enquanto há gente conectada.
  \item Pequenas diferenças de posição (menos de 2 segundos) são ignoradas
        de propósito, para evitar micro-ajustes constantes por causa da
        rede.
\end{itemize}

\section{Solução de problemas}

\begin{description}[style=nextline]
  \item[Não está sincronizando]
    Confira se o \textbf{Worker URL} e o \textbf{Código da sala} são
    exatamente iguais nas duas instalações da extensão (Passo 4).

  \item[A extensão não conecta]
    Confirme se o backend ainda está publicado --- rode
    \texttt{npx wrangler deploy} de novo dentro de \texttt{worker} caso
    necessário. Abra o console do navegador (\texttt{F12} → aba
    \textit{Console}) e procure mensagens que começam com
    \texttt{[YT Music Sync]} para ver o que está acontecendo.

  \item[A extensão sumiu no Firefox]
    Extensões temporárias somem quando o navegador reinicia --- repita o
    Passo 3 (Firefox).

  \item[Quero trocar de sala ou de backend]
    Basta abrir o popup da extensão de novo e alterar os campos --- a
    extensão reconecta automaticamente com os novos valores.
\end{description}

\end{document}
